\section{Discussion}
\label{sec:discussion}

\subsection{Cost accounting (indexing 1{,}000 documents)}

\begin{table}[h]
\centering
\caption{Indexing cost for 1{,}000 documents.}
\label{tab:cost}
\begin{tabular}{@{}lcc@{}}
\toprule
Method & Time & API cost \\
\midrule
\textbf{J-GraphRAG basic} (concepts + relations + embeddings) & \textbf{$\sim$5 min} ($\sim$1020 forwards) & \textbf{\$0} \\
\textbf{J-GraphRAG full} (+ roles + relation refinement) & \textbf{$\sim$8 min} ($\sim$2040 forwards) & \textbf{\$0} \\
Fast-GraphRAG (Qwen-7B generate) & $\sim$1000 min & \$0 (but unreliable) \\
Conventional GraphRAG (API) & $\sim$10 min & $\sim$\$5--15 \\
\bottomrule
\end{tabular}
\end{table}

\subsection{Methodological division of labor with interpretability research}
\label{sec:discussion:methodology}

The J-space paper proves the existence and causal privilege of the workspace
(interventional experiments: swap/steering/gradient pursuit); this work proves
the structure can be harvested by engineering. The key methodological
difference: \emph{corpus-grounded validation replaces interventional
validation}---in engineering extraction settings, ground truth in the data
(corpus validation, DF statistics, judge sampling) substitutes for causal
proof, cutting validation cost from months to minutes. We independently
replicated their Claude-scale findings on a 7B open model (position $-1$
readability, noise in the first third of layers, prefill modulation).

\subsection{Limitations}
\label{sec:discussion:limitations}

\begin{itemize}
\item Single-model validation (Qwen2.5-7B 4bit + fitted J-Lens); cross-model
generalization unknown.
\item Two domains (medical/novel); the L4 sample is small (14 questions), and
on novel L4 \emph{all} methods---including dense baselines---score near zero.
We currently suspect a question-set/judge artifact rather than a retrieval
failure, but this is not yet confirmed.
% TODO(Phase54): replace the following sentence with the confirmed diagnosis of
% the novel-L4 near-zero anomaly (question set vs. judge vs. retrieval).
\todo{Phase 54 diagnosis of novel-L4 near-zero pending.}
If the investigation confirms an evaluation artifact, the L4 numbers should be
re-measured on a corrected protocol and our retrieval conclusions stand; if it
instead reveals a genuine retrieval deficiency, the escape architecture of
\S\ref{sec:escape} predicts precisely this failure mode (queries linking to no
graph node) and the L4 results become a stress test for confidence gating.
\item The DeepSeek judge is non-deterministic ($\pm 0.03$--$0.05$);
borderline conclusions (novel 0.881) are limited by this.
\item Entity--entity relation edges for rare proper names are dominated by
co-occurrence default edges; scaling readout coverage is future work.
\end{itemize}

\subsection{Future work}
\label{sec:discussion:future}

\begin{enumerate}
\item \textbf{Larger L4 samples} to confirm the multi-hop gain (full novel
question set or the 2062-question subset).
% TODO(Phase55): fill in the confirmed multi-hop gain (or its absence) from the
% larger-L4-sample run; the transition below supports either outcome.
\todo{Phase 55 multi-hop gain on larger L4 sample pending.}
A confirmed gain would close the weakest link in the effect story and justify
graph expansion at L4; a null result would shift the multi-hop claim to
discrimination-only (\S\ref{sec:capability}) and redirect effort to generation
side.
\item \textbf{Cross-model generalization}: fitting J-Lenses for
Qwen3 / Gemma-3 / Llama-3.1 and verifying workspace readability---pretrained
lenses for these models already exist in the
neuronpedia/jacobian-lens ecosystem, making this a validation rather than a
training effort.
\item \textbf{Parameter count vs.\ lens accuracy ablation}: how readout
quality (concept accuracy, grounding AUC) scales with model size, and whether
larger models lift the structural ceiling on rare-entity readout
(\S\ref{sec:negative}).
\item \textbf{Escape architecture in practice}: measured confidence-gating
trigger rates and the true cost ratio.
\item \textbf{Qwen as dense encoder} to remove the last bge dependency on the
query side (optional; motivation is deployment simplicity).
\item \textbf{Tensor route at scale}: retest directional capability once
edge/concept counts grow.
\item \textbf{Multimodal extension}: applying workspace readout to
vision-language models so that figures, tables, and images can be indexed with
the same single-forward-pass extraction.
\end{enumerate}
